Work first conditionally on \(A\) and \((f_i)_{i\in\mathcal V_n}\), and define the proof-intermediate function
\[
 g_i(x):=\frac{f_i(1,x)+f_i(0,x)}{2},\qquad x\in[0,1].
\]
By \(\texttt{ass:bernoulli-design}\), conditional on \((U_i,f_i)_{i\in\mathcal V_n}\) and \(A\), every assignment vector has probability \(2^{-n}\). This probability does not depend on the latent types. The conditional-expectation identity
\[
 \Pr\!\left(Z=z\mid A,(f_i)_{i\in\mathcal V_n}\right)
 =\mathbb E\!\left[
 \Pr\!\left(Z=z\mid A,(U_i,f_i)_{i\in\mathcal V_n}\right)
 \mathrel{\big|}A,(f_i)_{i\in\mathcal V_n}
 \right]
 =2^{-n}
\]
therefore shows that the assignment remains product Bernoulli conditional on \(A\) and the response functions.

Fix \(i\in\mathcal V_n\) with \(N_i=m>0\). Because \(A_{ii}=0\), \(M_i\) is a function only of \((Z_j)_{j\ne i}\). Hence \(\texttt{ass:bernoulli-design}\) gives
\[
 M_i\sim {\rm Binomial}(m,1/2)
 \quad\text{and}\quad
 Z_i\ \text{is independent of}\ M_i
\]
under the conditional law just obtained. Define the unit-level contribution
\[
 \beta_i:=4\,\mathbb E_Z\!\left[
 Y_i(Z)(M_i-m/2)\mid A,(f_j)_{j\in\mathcal V_n}
 \right].
\]
By \(\texttt{ass:anonymous-interference}\), \(Y_i(Z)=f_i(Z_i,1/2+V_i)\) for this nonisolated unit. Averaging first over the independent variable \(Z_i\) replaces this outcome by \(g_i(1/2+V_i)\), while the definition of \(V_i\) gives \(M_i-m/2=mV_i\). Consequently,
\[
 \beta_i
 =4m\,\mathbb E_Z\!\left[
 g_i(1/2+V_i)V_i\mid A,(f_j)_{j\in\mathcal V_n}
 \right].
\]
If \(N_i=0\), then \(M_i-N_i/2=0\), so set \(\beta_i=0\). The construction in \(\texttt{def:li-wager-target}\) now yields
\[
 \bar\tau_{\mathrm{IND},n}=\frac1n\sum_{i\in\mathcal V_n}\beta_i.
\]

For the fixed nonisolated unit, let \(X_1,\ldots,X_m\) be independent random variables taking the values \(-1/2\) and \(1/2\) with equal probability, and put \(S_m=\sum_{j=1}^mX_j\). The preceding binomial law shows that \(V_i\) has the law of \(S_m/m\). Centering and independence give
\[
 \mathbb E S_m^2=\sum_{j=1}^m\mathbb E X_j^2=\frac m4,
 \qquad
 \mathbb E S_m^3=0.
\]
In the fourth-power expansion, the only terms with nonzero expectation are the one-index terms and the terms consisting of two pairs. Since \(\mathbb E X_j^2=1/4\) and \(\mathbb E X_j^4=1/16\),
\[
 \mathbb E S_m^4
 =\sum_{j=1}^m\mathbb E X_j^4
 +6\sum_{1\le j<k\le m}\mathbb E X_j^2\mathbb E X_k^2
 =\frac{3m^2-2m}{16}.
\]
Thus
\[
 \mathbb E V_i=0,
 \qquad
 \mathbb E V_i^2=\frac1{4m},
 \qquad
 \mathbb E V_i^3=0,
 \qquad
 \mathbb E V_i^4=\frac{3m^2-2m}{16m^4}\le\frac3{16m^2}.
\]

By \(\texttt{ass:third-derivative-bound}\) and the definition of \(g_i\),
\[
 \sup_{x\in[0,1]}\left|g_i^{(3)}(x)\right|\le L.
\]
Repeated use of the fundamental theorem of calculus on the segment joining \(1/2\) and \(1/2+v\), which is contained in \([0,1]\) for \(v\in[-1/2,1/2]\), gives
\[
 g_i(1/2+v)
 =g_i(1/2)+g_i'(1/2)v+\frac12g_i''(1/2)v^2+R_i(v),
 \qquad
 |R_i(v)|\le\frac L6|v|^3.
\]
Multiplying by \(v\), taking the conditional expectation, and using the four moment identities above makes the constant and quadratic Taylor terms vanish and gives \(4m\mathbb E V_i^2=1\). Therefore
\[
 \beta_i-g_i'(1/2)
 =4m\,\mathbb E\!\left[V_iR_i(V_i)\mid A,(f_j)_{j\in\mathcal V_n}\right],
\]
and hence
\[
 \left|\beta_i-g_i'(1/2)\right|
 \le \frac{4mL}{6}\mathbb E V_i^4
 \le \frac{4mL}{6}\frac3{16m^2}
 =\frac{L}{8m}.
\]
By \(\texttt{def:target}\), the contribution of unit \(i\) to \(\theta_n\) is \(g_i'(1/2)\). For an isolated unit, \(\texttt{ass:response-envelope}\) gives
\[
 |g_i'(1/2)|
 \le\frac{|\partial_xf_i(1,1/2)|+|\partial_xf_i(0,1/2)|}{2}
 \le B.
\]
Separating the nonisolated and isolated vertices and applying the triangle inequality proves
\[
\begin{aligned}
 \left|\bar\tau_{\mathrm{IND},n}-\theta_n\right|
 &\le \frac1n\sum_{i:N_i>0}\left|\beta_i-g_i'(1/2)\right|
 +\frac1n\sum_{i:N_i=0}|g_i'(1/2)|\\
 &\le \frac{L}{8n}\sum_{i:N_i>0}\frac1{N_i}
 +\frac Bn\sum_{i\in\mathcal V_n}\mathbf1\{N_i=0\}.
\end{aligned}
\]

We next derive an auxiliary bound involving only the fixed lower-order envelope. By \(\texttt{ass:response-envelope}\), \(\sup_x|g_i''(x)|\le B\). The first-order Taylor identity obtained from the fundamental theorem of calculus is
\[
 g_i(1/2+v)=g_i(1/2)+g_i'(1/2)v+T_i(v),
 \qquad
 |T_i(v)|\le\frac B2v^2.
\]
The same moment cancellation as above, followed by \(\mathbb E|V_i|^3\le(\mathbb EV_i^4)^{3/4}\), gives
\[
\begin{aligned}
 |\beta_i-g_i'(1/2)|
 &\le 2mB\,\mathbb E|V_i|^3
 \le 2mB\,(\mathbb E V_i^4)^{3/4}\\
 &\le 2mB\left(\frac3{16m^2}\right)^{3/4}
 =\frac{3^{3/4}B}{4\sqrt m}.
\end{aligned}
\]

Consider now any maintained triangular-array sequence with fixed \(B\) and \(\kappa\) and with
\[
 \liminf_{n\to\infty}\frac{\log\rho_n}{\log n}>-\frac12.
\]
Choose \(\varepsilon>0\) smaller than the gap between this lower limit and \(-1/2\). Then, for all sufficiently large \(n\),
\[
 \rho_n\ge n^{-1/2+\varepsilon}.
\]
Fix a vertex \(i\). By the independent uniform latent types in \(\texttt{ass:unit-sampling}\) and conditional edge independence in \(\texttt{ass:graphon-draw}\), after conditioning only on \(U_i=u\), the pairs consisting of \(U_j\) and the edge draw \(A_{ij}\), for \(j\ne i\), are independent and identically distributed. Marginalizing each \(U_j\) therefore gives
\[
 N_i\mid U_i=u\sim {\rm Binomial}(n-1,p_n(u)),
 \qquad
 p_n(u):=\rho_n\int_0^1W(u,v)\,dv.
\]
The conditions \(\texttt{ass:kernel-bound}\) and \(\texttt{ass:edge-probability}\) imply \(0\le p_n(u)\le\rho_n\kappa\le1\). Moreover, \(\texttt{ass:degree-regularity}\) gives, for almost every \(u\),
\[
 \mu_i(u):=\mathbb E[N_i\mid U_i=u]
 =(n-1)p_n(u)
 \ge\frac{(n-1)\rho_n}{\kappa}.
\]
Because every \(U_i\) is uniform by \(\texttt{ass:unit-sampling}\), this almost-everywhere inequality applies to all finitely many realized types with probability one at each \(n\).

For completeness, derive the required binomial lower-tail estimate. If \(S\sim{\rm Binomial}(m,p)\), \(\mu=mp\), and \(t=\log2\), then on \(\{S\le\mu/2\}\) one has \(e^{-tS}\ge e^{-t\mu/2}\). Markov's inequality and \(1-y\le e^{-y}\) consequently give
\[
\begin{aligned}
 \Pr(S\le\mu/2)
 &\le e^{t\mu/2}\mathbb E e^{-tS}
 =e^{t\mu/2}(1-p+pe^{-t})^m\\
 &\le\exp\!\left\{\frac{t\mu}{2}-\mu(1-e^{-t})\right\}
 =\exp\!\left\{-\frac{1-\log2}{2}\mu\right\}
 \le e^{-\mu/8},
\end{aligned}
\]
where the last step uses \((1-\log2)/2\ge1/8\). Since \(\mu_i(u)\ge(n-1)\rho_n/\kappa\), the event
\[
 \left\{N_i<\frac{(n-1)\rho_n}{2\kappa}\right\}
\]
is contained in \(\{N_i<\mu_i(u)/2\}\). Applying the preceding estimate conditionally on \(U_i\), integrating over \(U_i\), and taking a union bound yield, for
\[
 \mathcal E_n:=\left\{\min_{i\in\mathcal V_n}N_i
 \ge\frac{(n-1)\rho_n}{2\kappa}\right\},
\]
the probability bound
\[
 \Pr(\mathcal E_n^c)
 \le n\exp\!\left\{-\frac{(n-1)\rho_n}{8\kappa}\right\}
 \longrightarrow0.
\]
Indeed, the exponent diverges at least as a fixed multiple of \(n^{1/2+\varepsilon}\), which dominates \(\log n\).

On \(\mathcal E_n\) there are no isolated vertices and
\[
 \frac1n\sum_{i\in\mathcal V_n}\frac1{\sqrt{N_i}}
 \le\left(\frac{2\kappa}{(n-1)\rho_n}\right)^{1/2}.
\]
The auxiliary response-envelope bound thus implies
\[
 \frac{\left|\bar\tau_{\mathrm{IND},n}-\theta_n\right|}{\sqrt{\rho_n}}
 \le\frac{3^{3/4}B\sqrt{2\kappa}}{4\sqrt{n-1}\,\rho_n}
 \le\frac{3^{3/4}B\sqrt{2\kappa}}{4\sqrt{1-1/n}\,n^{\varepsilon}}
 \longrightarrow0
 \qquad\text{on }\mathcal E_n.
\]
Together with \(\Pr(\mathcal E_n^c)\to0\), this proves, independently of the additional Li--Wager source conditions,
\[
 \bar\tau_{\mathrm{IND},n}-\theta_n=o_P(\sqrt{\rho_n}).
\]

Finally restrict to the further intersection with every source condition recorded in \(\texttt{lem:li-wager-theorem-six-regime}\), specialize the source treatment probability to one half as required by \(\texttt{ass:bernoulli-design}\), and suppose the equation~(43) normalization is well defined. The cited lemma then supplies the published standardized Gaussian limit centered at \(\bar\tau_{\mathrm{IND},n}\), with center measured on the \(\sqrt{\rho_n}\) scale. The last display says that replacing that center by \(\theta_n\) adds a term converging to zero in probability in the same normalization. To verify the needed additive Slutsky step directly, write the published standardized statistic as \(T_n\), the normalized center change as \(R_n=o_P(1)\), and let \(\Phi\) be the standard-normal distribution function. For every real \(x\) and every \(\eta>0\),
\[
 \Pr(T_n\le x-\eta)-\Pr(|R_n|>\eta)
 \le \Pr(T_n+R_n\le x)
 \le \Pr(T_n\le x+\eta)+\Pr(|R_n|>\eta).
\]
Since the cited limit gives \(T_n\Rightarrow N(0,1)\) and \(R_n=o_P(1)\), taking lower and upper limits and then letting \(\eta\downarrow0\) uses continuity of \(\Phi\) to give \(T_n+R_n\Rightarrow N(0,1)\). Thus the published limits transfer to centering at \(\theta_n\) only on the stated source-class intersection. The nonzero degree-exact remainder proved above also confirms that the two finite-sample targets have not been identified with one another by definition.